%=====================================================================
%  7. Conclusion
%=====================================================================
\section{Conclusion}
\label{sec:conclusion}

We presented \sysname{}, a multi-agent framework for personalized domain
knowledge generation in vertical skill training. Its distinctive feature is
structural rather than algorithmic: verification is placed between generation
and delivery, the difficulty rule is explicit and independently recomputable,
and adaptation is expressed as a published decision table rather than left
implicit in model weights. Across \todo{61} end-to-end cases spanning
\todo{14} topics and \todo{12} learner profiles, the full configuration
outperforms single-pass and ablated alternatives on adaptation accuracy,
coverage, and hallucination rate, with the debate stage and retrieval
grounding contributing the largest measurable shares.

We are deliberately explicit about what this paper does \emph{not} establish:
we demonstrate artifact-level quality and profile fit, not improved learning
outcomes. Closing that gap requires a controlled human study, which is the
central item on our agenda, together with replication on a second model family
and a longitudinal deployment that tracks profile evolution.

More broadly, the results suggest that for educational generation, the
bottleneck is no longer whether a model can produce plausible material, but
whether a system can demonstrate that the material is correct for \emph{this}
learner. Architectures that make verification and adaptation inspectable
appear better suited to that requirement than architectures that pursue
quality purely through scale.
